\section{The quantitative Weil dictionary on the fixed choice}
\label{sec:quantitative-weil-dictionary}

This section strengthens WA11--WA16 and CW23--CW25 in a different
direction: it proves the all-conjugate scalar and local-unit assertions,
and the exact denominator and coefficient bounds in the source's
\texttt{DictAt}.  The input is the \emph{same supplied} arithmetic host,
\texttt{Choice} $\mathfrak c$ and proof $\mathfrak c.\mathrm{Is}$.  In
the literal source endpoints we retain their arguments
$\texttt{FunPr.Hypotheses}_2$ and $\forall D,\texttt{NormSum.Inp}(D)$;
we do not assert the construction of those whole packages here.
The substantive ingredients of these endpoints are proved below,
using the actual Weil action and the already reconstructed theta and
Fock comparisons.  In particular there is no independent assumption
of the desired scalar, dictionary, or simultaneous local-unit property.

Write $E/F$ for the CM extension, $f=[F:\mathbb Q]=D.R.f$, and
$p$ for a literal \texttt{DefPair}.  Thus $p$ is either a top edge or
a rightmost $\chi$ edge; the $\mu$ and core cases are excluded.
Let $a,b$ be the ranks of its oriented source and target.  For finite
points $x,y$, let $L_s(x,y),L_t(x,y)$ be the source's
\texttt{srcLat}, \texttt{tgtLat}, in its fixed rational orthonormal
frames.  For a $\chi$ edge the vertex order is reversed and the first
Hermitian space is negated, exactly as in \texttt{toPair}.
The involution $\operatorname{opp}_p$ is complex conjugation on that
edge and the identity on a top edge.  Put
\[
 A_0=\operatorname{La}_0(D,\operatorname{tor}_0(D,\mathfrak c_K,\delta)),
 \qquad A_{x,y}=\operatorname{La}(L_s(x,y),L_t(x,y)),
 \qquad G_{x,y}=I_{A_0,A_{x,y}}e_0.                         \tag{QD.1}
\]
Here $I$ is the source's \emph{positively normalized} finite
Heisenberg intertwiner, $e_0=\mathfrak c.S.Z.e_0(p)$,
$\delta=\mathfrak c.S.Z.\delta$, and $c_0=\mathfrak c.S.Z.c_0(p)$.
Neither $e_0$ nor $c_0$ is rescaled.  In particular $c_0\ne0$ need
not have absolute value one.  Denote by $J_{x,y}(\alpha)$ the actual
Fock monomial jet of
$\omega_p(\operatorname{toPair}_p(xu_s,yu_t))v_p$, and put
$P_\alpha=\prod_\nu\operatorname{period}(\delta,\nu)^{\alpha_\nu}$.
The statement \texttt{JetKap} is precisely
\[
 J_{x,y}(\alpha)=c_0\operatorname{opp}_p
       \left(P_\alpha\,\partial^\alpha(\kappa G_{x,y})(0)\right)
                         \quad\hbox{for every }\alpha.     \tag{QD.2}
\]
The notation $\partial^\alpha$ means the source's
\texttt{applyAt0} on the monomial; it includes the multi-factorial
relative to the Taylor coefficient.

\paragraph{QD1: extracting arithmetic from a positive power.}
Suppose $z^m=u$, $m>0$, and $u\in\overline{\mathbb Q}$.
Then $z\in\overline{\mathbb Q}$.  If every complex embedding of
$u$ has modulus one, so does every complex embedding of $z$.
For a fixed rational prime $q$, if every embedding of $u$ into
$\mathbb C_q$ has norm one, the same is true of $z$.
Indeed $X^m-u$ proves algebraicity over the algebraic numbers,
and under any of these embeddings the equality gives
$|z|^m=|u|=1$.  The nonnegative real number $|z|$ is therefore
one.  Products, inverses, and complex conjugates preserve both
assertions, with the latter interpreted by precomposition with
the conjugation automorphism of $\overline{\mathbb Q}$.
Roots of unity satisfy all of them and are algebraic integers.
We use $\mathrm{Mod}_1(z)$ for the all-complex-embedding assertion,
including algebraicity, and $\mathrm{Unit}_q(z)$ for the assertion
at every $\mathbb C_q$ embedding.  These are the source's
\texttt{Mod1All} and \texttt{UnitAt}.  A single complex norm
identity is not an alternative proof of either statement.

\paragraph{QD2: finite Hecke values have the required conjugate norms.}
Let $\chi$ be either actual splitting Hecke character and let $z$
be a finite idele, with infinite component one.  Then
$\mathrm{Mod}_1(\chi(z))$.  Moreover,
\[
 z_P=1\quad\hbox{for every }P\mid q
                  \quad\Longrightarrow\quad
                       \mathrm{Unit}_q(\chi(z)).          \tag{QD.3}
\]
Here and below a circle value is included into $\mathbb C$.

The supplied \texttt{ResEta} may be tested against HF9's actual
quadratic character.  Its norm triviality makes $\chi$ trivial
on positive real scalars at infinity.  On each complex factor,
polar decomposition and the elementary classification of circle
characters give an integer $m_\tau$; its infinity value is
$\prod_\tau(z_\tau/|z_\tau|)^{m_\tau}$.  One way to prove that
classification is to lift the continuous map along
$t\mapsto e^{2\pi i t}$ near zero, extend using the homomorphism
law, and observe that its additive lift has integer value at $1$.

The finite integral idelic units form a profinite compact group.
A sufficiently small open subgroup maps into an arc of the
circle containing no nontrivial subgroup, so it maps to one.
The character on the whole compact group consequently has finite
image.  Finite ideal-class number now gives $h>0$, $b\in E^\times$
and an integral finite unit idele $w$ with $z^h$ equal, at finite
places, to $w b$.  If $z_P=1$, taking valuations gives $v_P(b)=0$.
Raising once more to the order $M$ of $\chi(w)$, and using
triviality on principal ideles, gives
\[
 \chi(z)^{hM}\prod_\tau
       \left({\tau b\over|\tau b|}\right)^{M m_\tau}=1,
 \qquad
 \chi(z)^{2hM}=
       \prod_\tau\left({\overline{\tau b}\over\tau b}
                  \right)^{M m_\tau}.                    \tag{QD.4}
\]
Every embedding of the CM field sends its involution to complex
conjugation.  Thus every complex conjugate of the right side has
modulus one.  If $v_P(b)=0$ at all $P\mid q$, both $b$ and $b^{-1}$
are integral over $\mathbb Z_{(q)}$.  Explicitly each has an integer
denominator prime to $q$, obtained by clearing its finitely many
bad primes outside $q$.  Their images in every $\mathbb C_q$ have
norm at most one, hence exactly one.  This applies to all the
factors in (QD.4).  QD1 proves the assertions.  No continuity of
an arbitrary field embedding, and no universal uniqueness theorem
for \texttt{IsEta}, has been used.

\paragraph{QD3: rational unitary determinants and exact local lattices.}
For $U\in U_n(E)$ in the fixed standard Hermitian coordinates,
$\mathrm{Mod}_1(\tau\det U)$ for every chosen embedding $\tau$.
If full $\mathcal O_E$-lattices $M,UM$ are both
\texttt{EqAt} $q$ to $\mathcal O_E^n$, then also
$\mathrm{Unit}_q(\tau\det U)$.

The first assertion follows by taking determinants in $U^*U=1$
and then applying each embedding of the CM field.  For the
second, \texttt{EqAt} supplies integer multipliers prime to $q$
for the inclusions in both directions.  Combining the inclusions
gives integers $s,t$ prime to $q$ for which $sU$ and $tU^{-1}$
have integral entries.  Under every embedding into $\mathbb C_q$,
their determinants have norm at most one, whereas $s,t$ have
norm one.  Hence both $\det U$ and its inverse have norm at most
one.  This proves equality.  The argument needs the two actual
integral lattices, not just equality of their dimensions or their
rational spans.  Conversely, equality after localization at all
$P\mid q$ gives the source's integer-denominator witnesses:
clear the finitely many other prime denominators on generators
in both directions by integers prime to $q$.

\paragraph{QD4: a finite quadratic Gauss sum has finite-order phase.}
Let $B$ be a finite abelian group and $q:B\to\mathbb C^\times$
take root-of-unity values, with $q(0)=1$ and
$b(x,y)=q(x+y)/(q(x)q(y))$ a bicharacter.  Then
$S=\sum_{x\in B}q(x)$ is either zero or a positive real number
times a root of unity.  This holds at $2$ as well as at odd primes.

On the radical $R$ of $b$, $q$ is a character.  If that character
is nontrivial, translation by an element of $R$ proves $S=0$.
Otherwise $q$ descends to $B/R$ and $S=|R|\overline S$.
It is enough to consider the nondegenerate case, of order $N$.
On $\mathbb C^B$ put
\[
 F_{x,y}=N^{-1/2}b(x,y),\qquad T_{x,y}=1_{x=y}q(x).
\]
Character orthogonality gives $F^2f(x)=f(-x)$ and $F^4=1$.
Translation of the sum gives
$\sum_yq(y)b(y,t)=S/q(t)$, so direct matrix multiplication gives
\[
 FTF={S\over\sqrt N}\,T^{-1}F^{-1}T^{-1},
              \qquad (FT)^3={S\over\sqrt N}\,1.           \tag{QD.5}
\]
The matrices $F,T$ are unitary.  Thus $|S|=\sqrt N$.
Their determinants are roots of unity, since $F^4=1$ and the
diagonal entries of $T$ have finite order.  Taking determinants
in (QD.5) shows that $(S/\sqrt N)^N$, and hence $S/\sqrt N$,
is a root of unity.  A nontrivial constant $q(0)$, when present,
can first be divided out.  This proof needs no prime-power
classification or choice of a square root of an algebraic phase.

\paragraph{QD5: phase control for the actual positive intertwiners.}
For three commensurable self-dual period lattices with their
actual compatible semicharacters, the scalar in composing their
positively normalized intertwiners is a root of unity.  Only this
conclusion, rather than a specified order dividing eight, is
needed below.

Choose a common integral lattice $D$ small enough that all three
lattices lie in $D^\#$ and all their semicharacters agree on
$D$.  Multiplication by an even integer achieves the last
condition.  The calculation is then in the finite Heisenberg
module $D^\#/D$.  The normalized intertwiner is a positive
normalization of the sum of the actual operators
$\alpha(x)U(x)$ over the target lattice modulo the intersection.
Their covariance identifies its image with the target theta line;
VI5 proves it is an isometry and is nonzero.  On that line the
composition therefore differs by a scalar of absolute value one.

For completeness, its phase is a finite quadratic Gauss phase.
Expand the trace of the three finite lattice projectors.
Heisenberg orthogonality kills every summand except
$x+y+z=0$.  On this kernel the remaining summand is
\[
 \alpha_A(x)\alpha_B(y)\alpha_C(z)
          e(x,y)e(x+y,z),                                \tag{QD.6}
\]
where $e(v,w)$ is the phase in the \emph{actual} law
$U(v)U(w)=e(v,w)U(v+w)$, with the same convention as VI2
and CA1.  All its values are roots of unity: the vectors are
rational modulo a full integral lattice and their denominators
are bounded.  The semicharacter law shows directly that the
polarization of (QD.6) is a bicharacter.  The omitted trace,
index and normalization factors are positive real numbers.
QD4 therefore proves that the nonzero composition scalar has
finite-order phase.  Its absolute value is one, so it is itself
a root of unity.  Passing through the finite quotient is legitimate
because every projector sum and every trace in this calculation
is finite.  This supplies the phase assertion used here without
assuming equality of positive and signed intertwiners in odd rank.

\paragraph{QD6: the determinant controls a change of theta product basis.}
Let $b=(b_i)$ and $b'=(b'_i)$ be rational orthonormal bases of
the same positive Hermitian $K$-space used in the source's
\texttt{pullB}, and put
$G_{ij}=\phi_0(h(b'_j,b_i))$.  For the same normalized one-variable
theta function, period and un-killed half-point, write $\theta_b$
and $\theta_{b'}$ for the two product theta functions.  There are
$c\ne0$ and a root of unity $\zeta$ such that
\[
 \theta_{b'}=c I_{A_b,A_{b'}}\theta_b,
                         \qquad c^2=\zeta\det G.          \tag{QD.7}
\]

The theta-line theorem gives the first identity.  Both products
have the same positive mean square, the product of the fixed
one-variable mean squares; VI5's isometry therefore gives
$|c|=1$.  First suppose $d=\det(1-G^*)\ne0$.  The coordinate
unitary map $u$ has $\det(1-u)=d$.  Take a common sublattice
$\Lambda'\subset\Lambda$ on which $u\Lambda'\subset\Lambda$
and the two semicharacters agree.  Expand the reproducing theta
kernel, and unfold its integral on a fundamental domain of
$\Lambda'$ by the cosets of $(1-u)\Lambda'$ in $\Lambda$.
For $B=(1-u)^{-1}$ the term belonging to $x$ has integral
\[
 {C_H\over\det(1-u)}\,
       \alpha(x)\exp\{-\pi i\,\operatorname{Im}H(Bx,x)\},
                         \qquad C_H>0.                   \tag{QD.8}
\]
Here $C_H$ is the positive Haar scale of the Gaussian integral.
To verify (QD.8), translate the Gaussian by $Bx$ and complete
the square.  Its real quadratic part is positive definite since
$\operatorname{Re}H((1-u)z,z)=\|z-uz\|_H^2/2$.
Unitary diagonalization reduces the remaining integral to
$\int_{\mathbb C}e^{-\pi(1-\lambda)|z|^2}dz=(1-\lambda)^{-1}$.
Gaussian domination justifies both integration and unfolding.
The remaining finite sum is quadratic: its polarization follows
from
$\operatorname{Im}H(Bx,y)-\operatorname{Im}H(By,x)
=\operatorname{Im}H(x,y)$.  QD4 gives finite-order phase.

There is also an exact way to read this same integral on the
theta line.  Substitute the first identity of (QD.7) into the
kernel.  Orthogonality on the common fundamental domain kills
all nontrivial translate cosets; the zero coset contributes its
positive volume.  Thus $cd$ has finite-order phase.  Unitarity
of $G$ gives
\[
       \overline d=(-1)^r\det(G)d,
       \qquad |d|^2=d\overline d.                         \tag{QD.9}
\]
If $\omega=cd/|d|$ is its root-of-unity phase, (QD.9) gives
$c^2=(-1)^r\omega^2\det G$, proving (QD.7).

Finally choose a rational norm-one $z\in K$ avoiding the finitely
many forbidden eigenvalues, and insert $b''=zb$.
The values $(t+\sqrt{-d})/(t-\sqrt{-d})$, $t\in\mathbb Q$,
give infinitely many choices.  Both consecutive changes now
satisfy the determinant condition.  Compose their identities
and use QD5; their determinants multiply to $\det G$ and all
extra composition scalars are roots of unity.  This proves
(QD.7) with no fixed-vector exception.  The zero-dimensional
case has $c=\zeta=1$.

\paragraph{QD7: rational changes of the period lattice preserve arithmetic phase.}
For rational unitary $U\in U_a(E)$ and $U'\in U_b(E)$ and
self-dual $M,M'$, there is $\epsilon'$ with
$\mathrm{Mod}_1(\epsilon')$ such that
\[
 I_{A_0,\operatorname{La}(M,M')}e_0
          \bigl(\operatorname{tPull}(U,U')z\bigr)
 =\epsilon' I_{A_0,\operatorname{La}(UM,U'M')}e_0(z).      \tag{QD.10}
\]
If $M,M',UM,U'M'$ are all \texttt{EqAt} $q$ to their standard
lattices, then $\mathrm{Unit}_q(\epsilon')$.

Use the source's product-theta presentation of this same $e_0$,
$e_0=\kappa_b I_{A_b,A_0}\theta_b$.  Apply the linear isometry
$X\mapsto U'XU^*$ to $b$ and the lattice.  Finite sums show
that pullback commutes with each intertwiner; the quotient index
and its positive square root are unchanged.  QD6 supplies the
scalar for $b'_i=U'b_iU^*$, and QD5 supplies the three composition
phases.  The prefactor $\kappa_b$ occurs on both sides and cancels.
In particular $\epsilon'$ is the product of that $c$ with roots
of unity and their inverses, exactly the source's
$c\mu_1\mu_3\mu_2^{-1}$ (where $\mu_2^{-1}=\mu_2^7$
when its original eighth-power assertion is used).

The determinant occurring in (QD.7) is
\[
 \phi_0\det\bigl(h(U'b_jU^*,b_i)\bigr)
   =\prod_{\tau:E\hookrightarrow\mathbb C,\ \tau|_K=\phi_0}
          \tau(\det U')^a\tau(\det U^*)^b.                \tag{QD.11}
\]
This follows by vectorizing $X\mapsto U'XU^*$ over $E$ and
taking its determinant after restriction of scalars to $K$;
an orthonormal basis introduces no further Gram determinant.
QD3 controls (QD.11) under every required embedding, and QD1
extracts its square root.  This proves both assertions about
$\epsilon'$.  No evenness of each rank separately is used;
in particular CA5's more restrictive identification of signed
and positive intertwiners has not been invoked.

\paragraph{QD8: real vacuum characters and rational finite moves.}
For the actual supplied Weil action on definite spaces, the
scalar of its real compact group on a fixed nonzero Fock/lattice
vacuum is an integer determinant character.  In orthonormal
coordinates it has the form
\[
 c(g,h)=\prod_\tau\det(g_\tau)^{m_\tau}
                         \det(h_\tau)^{n_\tau},
                         \qquad m_\tau,n_\tau\in\mathbb Z. \tag{QD.12}
\]
The exponents are the winding numbers of the \emph{given}
action; we do not substitute a new action or character for it.

The real compact action preserves the Fock equations, commutes
with finite Heisenberg operators and preserves the lattice
conditions.  Theta-line uniqueness therefore gives a scalar on
this line.  Its Haar norm and continuity make it a continuous
circle-valued character.  Restrict a factor $U(n)$ to the diagonal
torus.  The circle calculation in QD2 gives integer exponents.
Conjugation by permutation matrices makes them all equal.
Every unitary matrix is unitarily diagonalizable and a scalar
character is conjugation invariant, so the character is the
corresponding power of the determinant.  Apply this to each
factor; $U(0)$ contributes one.  This proof applies also to the
negative definite source of a $\chi$ edge.  It avoids the source
route's auxiliary universal \texttt{IsEtaTilde} uniqueness issue.

Let $r,r'$ be the finite components of rational unitary elements,
with their actual coordinate matrices $U,U'$.  Rational automorphy
of the same Weil action and its finite/real commutation give the
finite action as rational pullback followed by the inverse real
vacuum character.  Differentiating this identity in the fixed
frames gives the source's \texttt{JetTh} substitution
\[
 \Theta(z)\longmapsto
       \epsilon\,\Theta(\operatorname{tPull}(U_s,U_t)z).
                                                               \tag{QD.13}
\]
Here $(U_s,U_t)=(U,U')$ for a top edge and $(U',U)$ for a
$\chi$ edge, with its stipulated negation and conjugation.
If $c$ is the scalar before the finite/real rearrangement,
$\epsilon=c^{-1}$ in the top case and $\epsilon=c$ in the
$\chi$ case, since $\overline c^{-1}=c$.  Formula (QD.12) and
QD3 prove $\mathrm{Mod}_1(\epsilon)$ and
$\mathrm{Unit}_q(\epsilon)$ when the four before/after vertex
lattices are standard at $q$.  The actual periods and $c_0$
are unchanged.  Applying QD7 to (QD.13) proves the source's
rational propagation of \texttt{JetKap}, with extra scalar
$\epsilon\epsilon'$ and precisely these four lattice premises.

\paragraph{QD9: the standard-lattice starting scalar is a root of unity.}
Suppose both actual moved vertex lattices are globally standard
in the fixed frames.  Then (QD.2) holds with a root of unity
$\kappa_0$.  Indeed the finite elements $xu_s,yu_t$ preserve
the original full self-dual lattices.  Their stabilizers are
compact open subgroups: after choosing a full rational lattice
frame, this is the compact integral matrix group, with the
closed Hermitian equations imposed.  The same Fock/lattice
vacuum line is preserved.  Its scalar character is continuous
and has values in the circle.  A small open subgroup maps to
an arc containing no nontrivial subgroup; compactness makes
the resulting quotient finite.  The scalar of each element
is consequently a root of unity.

Equivalently one may use the actual smooth-vector principal
level: some positive powers of the two stabilizer elements lie
in it, so a common positive power fixes the nonzero vacuum.
For standard lattices $I_{A_0,A_0}e_0=e_0$.  The fixed
\texttt{kernelNorm\_constant\_spec} then gives (QD.2), with
$\kappa_0$ the scalar or its conjugate according to $p$.
QD1 gives all its complex and local norms.  This argument
does not change the period-normalized $e_0$ into a unit vector.

\paragraph{QD10: the distinguished scalar norm is exactly one.}
Any scalar satisfying (QD.2) obeys $|\kappa|=1$.
Here is the normalization calculation which will remove the
positive factor in a split move.  Put
$N_0=\|e_0\|_{A_0,\mathrm{mean}}^2>0$, with the exact mean-square
normalization of VI5.  The original unit vacuum and its moved
unit vacuum have Gaussian entire restrictions.  Their Fock jet
identities identify those restrictions with
$c_0 e_0$ and $c_0\kappa G_{x,y}$ after the same period scaling.
For a $\chi$ edge conjugate both identities first.

Partition a rational adelic fundamental domain by the finite
lattice cosets and apply its rational quasi-periodicity.  The
finite coset count multiplies both volume and integral, so it
cancels in the normalized mean.  The real restriction is
Gaussian times its entire factor, and the substitution by
$\delta$ changes the Hermitian Gaussian and covolume by the
same coordinate Jacobian.  Thus the Haar unit-norm assertion
is exactly the period-lattice mean-square assertion.  This
can also be checked on one common real fundamental domain:
translate all pieces there; the absolute square of each phase
is one, and the normalized averages agree.  VI5 gives
$\|G_{x,y}\|_{A_{x,y},\mathrm{mean}}^2=N_0$.  Consequently
\[
                   |c_0|^2N_0=1,
                   \qquad |c_0\kappa|^2N_0=1.             \tag{QD.14}
\]
Cancel the first positive quantity to obtain $|\kappa|^2=1$.
All derivatives determine the same entire function: if needed,
their Taylor series agree near zero, then on every complex line
by the identity theorem.  No equality of the two norms is
inferred from jets without this entire-function step.  Formula
(QD.14) is only a distinguished-complex assertion; QD2, QD3
and QD6 supply the arithmetic information independently.

\paragraph{QD11: reading a split move as the same lattice intertwiner.}
Let $e$ be a split-place idempotent and let $s,s'$ be actual
finite unitary elements supported on $e+\bar e$.  Choose their
literal determinant lifts $x_0,y_0$, supported on $e$ (or on
$\bar e$ after the $\chi$-edge conjugation).  If (QD.2) holds
at $x,y$, it holds at $sx,s'y$ with scalar
\[
       r\zeta\,\chi_T(x_0)^{-1}\chi_S(y_0)\kappa,
                 \qquad r>0,\quad \zeta\text{ a root of unity}.
                                                               \tag{QD.15}
\]
This is an identity for the actual action, not merely the
existence of an unrelated vector on the same complex line.

Here is the finite calculation behind it.  Let $\Lambda$ and
$\Lambda'$ be the old and new self-dual finite Hom lattices.
They agree outside $e+\bar e$.  Choose representatives $R,R'$
for their two changed half-lattices modulo their intersections
with the old ones.  The two halves are isotropic.  Character
orthogonality therefore makes
\[
                  F=\sum_{r\in R}\sum_{r'\in R'}
                                      U(r)U(r')v          \tag{QD.16}
\]
a nonzero vector fixed by the new lattice, with the unchanged
real Fock equations.  To check nonvanishing and phase, use the
coinvariant functional of WA5: it is one at $v$, invariant on
the $\bar e$ half, and takes value $1_{\Lambda}(r)$ on the
$e$ half.  The operators in (QD.16) commute on these two new
lattice halves, since their full symplectic pairings have
trivial character.  Choose zero as the representative of the
zero coset.  The functional then gives $|R'|>0$ on (QD.16).
The split clause of the given Weil action gives the positive
multiple of $\chi_T(x_0)^{-1}\chi_S(y_0)$ on this functional.
Vacuum-line uniqueness identifies $\omega(s,s')v$ with (QD.16)
with exactly that phase.  Its norm fixes only the positive
factor, which is kept explicit at this point.

To read (QD.16) in the period model, approximate each
$r+r'$ by a rational Hom vector $\gamma$ modulo
$2(\Lambda\cap\Lambda')$.  Such approximation follows from
the finite Chinese remainder theorem on the finitely many
denominator primes; the rational vector may be chosen in
$\Lambda'$ at every finite place.  The remaining double
lattice vector fixes the old vacuum.  The Heisenberg law
turns its coefficient into the actual semicharacter of the
period vector $P\gamma$.  There is no undetermined phase here:
apply the same calculation to the nonzero new vacuum; since
both half-operators fix it, its period semicharacter must be
the coefficient just obtained.  The pairs $(r,r')$ index
exactly the cosets of the new period lattice modulo its
intersection with the old one.  Injectivity follows by
projecting a putative equality onto the two halves; surjectivity
follows by decomposing a new vector into $e,\bar e$ and the
unchanged complement.  Thus the sum reads as the positive
intertwiner $I_{A_{x,y},A_{sx,s'y}}$, up to a positive real
factor.  Compose it with $I_{A_0,A_{x,y}}$ and use QD5.  This
proves (QD.15).  On a $\chi$ edge perform this calculation on
the conjugated positive model, then conjugate back; the two
characters and the reversed source/target order are exactly
those displayed in (QD.15).

\paragraph{QD12: split propagation is arithmetic at every avoided prime.}
In (QD.15), $r=1$.  Both before and after scalars satisfy QD10;
the root of unity and the two actual circle-character values
also have complex norm one.  Taking norms in (QD.15) therefore
gives $r=1$.

If the supporting base prime does not lie over a rational prime
$q$, the determinant lifts $x_0,y_0$ are one at every $P\mid q$.
QD2 shows that the extra scalar
\[
                       \zeta\chi_T(x_0)^{-1}\chi_S(y_0)
                                                               \tag{QD.17}
\]
has $\mathrm{Mod}_1$ and $\mathrm{Unit}_q$.  The involution on
idempotents preserves this property.  Inverses of a supported
unitary element have the same support, so the conclusion
propagates through any finite word of supported elements and
their inverses.  Here \texttt{splitOff} is the algebraic subgroup
generated by such elements, not a topological closure; no limit
of algebraic numbers is being declared algebraic.

\paragraph{QD13: norm-one corrections with controlled nonsplit denominators.}
Fix a finite set $S$ of base-field primes, a norm-one $u\in E$,
and any required finite precision.  There exists a rational
norm-one $u'$ approximating $u$ at $S$ and approximating $1$
to any prescribed precision at another finite disjoint set.
At every conjugation-fixed prime $P$, $v_P(u')=0$.

For the approximation assertion use the elementary Hilbert--90
formula $u=\alpha/\bar\alpha$.  It follows directly by taking
$\alpha=\theta+u\bar\theta$ for a choice of $\theta$ making
this expression nonzero; if all such expressions vanished,
conjugation would be the identity.  The split local case is the
same explicit ratio of two components.  At the finitely many
places choose nonzero local $\alpha_v$ giving the desired
norm-one value.  Field weak approximation gives one
$\alpha\in E$ sufficiently close to all of them; the map
$\alpha\mapsto\alpha/\bar\alpha$ is continuous there.
The valuation assertion follows from
$v_P(\bar\alpha)=v_P(\alpha)$, including at ramified primes.

This gives the source's level correction for a unitary
quasireflection.  Write
\[
 R_{v,u}=1+(u-1)P_v,\qquad
              P_v(w)=v\,{h(w,v)\over h(v,v)}              \tag{QD.18}
\]
with the linear-variable convention adjusted to the actual
Hermitian form.  Its inverse is $R_{v,u^{-1}}$.  The projector
$P_v$ has integral lattice entries outside finitely many primes.
Choose $u'$ close to $u$ at the specified level primes and close
to $1$ at every other nonsplit prime where $P_v$ has a denominator.
At the latter places both $(u'-1)P_v$ and
$(u'^{-1}-1)P_v$ are integral; at every other nonsplit place
this follows from $v_P(u')=0$.  Thus the corrected reflection
and its inverse preserve the lattice outside the level primes
and split primes.  Its difference from the original reflection
is as small at the level primes as prescribed.  The finite
local inclusions can be expressed with the source's global
integer multipliers prime to those primes, by clearing
denominators outside them.  This proves the support control
in \texttt{normOne\_approx} and in the reflection correction
on the actual definite spaces, without a prime-existence or
strong-approximation assumption.

\paragraph{QD14: rational, split and stabilizer factorization.}
Let $g$ be a finite unitary element of an actual definite vertex,
let $\Lambda$ be its fixed full lattice, and let $Q$ be a finite
set of rational primes.  There is an exact factorization
\[
                 g=r s k,                               \tag{QD.19}
\]
where $r$ is the finite component of a rational unitary element,
$s$ is a finite product of split-supported elements away from
$Q$, and $k$ stabilizes $\Lambda$ everywhere.  At primes over
$Q$, $k$ can satisfy any prescribed principal congruence.

First recall why rational unitary weak approximation at finitely
many finite places is available here.  Multiply the desired local
unitaries by one rational norm-one scalar avoiding their finitely
many forbidden eigenvalues.  On this open set the Cayley map
\[
      X\longmapsto(1-X)(1+X)^{-1},\qquad X^*=-X,           \tag{QD.20}
\]
is a rational chart for the unitary group.  Its inverse has the
same form.  The skew-Hermitian matrices form an $F$-vector space;
approximate its coordinates by elements of $F$.  The nonvanishing
determinant conditions remain open, and the Cayley map preserves
the required precision.  Undoing the scalar gives a rational
$\gamma_0$ with $\gamma_0 g$ arbitrarily close to one at the
chosen places.  This is just finite field weak approximation.

Choose those places to include all primes over $Q$ and every
place where $g$ or $g^{-1}$ fails to preserve $\Lambda$.
Factor $\gamma_0$ into finitely many unitary quasireflections,
as in WA10.  Apply QD13 to their norm-one parameters, with
enough precision that the product $\gamma$ still has
$\gamma g$ integral and invertible at the chosen places and
has the requested congruence over $Q$.  At other nonsplit
places every corrected factor and its inverse preserve
$\Lambda$.  Thus $\gamma g$ can fail to preserve $\Lambda$
only at a finite set of split base primes away from $Q$.

Let $s$ agree with $\gamma g$ at precisely those split primes,
and be one elsewhere; let $k$ be one there and agree with
$\gamma g$ elsewhere.  Each support cut is fixed by conjugation,
so both matrices are unitary; they commute and multiply to
$\gamma g$.  The matrix $k$ and its inverse preserve the
lattice at every place.  The rational lattice is the intersection
of its localizations, so $k$ is in the actual adelic stabilizer.
The element $s$ is the finite product of its single split-pair
pieces.  Taking $r=\gamma^{-1}$ proves (QD.19).
This argument proves the definite-vertex instance needed here;
it does not label every more general source approximation
theorem as reconstructed.

\paragraph{QD15: propagating one scalar while protecting a finite prime set.}
Fix finite $Q$ such that both original vertex lattices at $x,y$
are \texttt{EqAt} standard at every $q\in Q$.  There exists
$\kappa_Q$ satisfying (QD.2), with $\mathrm{Mod}_1(\kappa_Q)$
and $\mathrm{Unit}_q(\kappa_Q)$ for all $q\in Q$.

Apply QD14 separately to $xu_s$ and $yu_t$ and write
$x=r s x_0$, $y=r's'y_0$, with
$x_0=k u_s^{-1}$ and $y_0=k'u_t^{-1}$.
The actual fixed-frame lattices of $x_0,y_0$ are globally standard.
QD9 supplies their root-of-unity scalar.  Apply QD12 to the
finite words $s,s'$, using one factor on one side at a time.
All intervening words act trivially at $Q$, so the two resulting
lattices remain standard there.  All scalars acquired thus far
have the required units at $Q$.  Finally apply QD8 and QD7
to $r,r'$.  The before lattices are standard at $Q$ by construction,
and the after lattices are the original ones, standard at $Q$
by hypothesis.  The four lattice premises of QD3 are therefore
present.  The product of the scalars proves the assertion.
All operations take place on the fixed original frames and on
the exact actual lattice $A_{x,y}$.

\paragraph{QD16: uniqueness gives the simultaneous all-prime scalar.}
The scalar in (QD.2) is unique.  The function $G_{x,y}$ is entire
and nonzero by theta-line uniqueness and VI5's invertibility of
$I_{A_0,A_{x,y}}$.  Hence at least one Taylor coefficient is
nonzero: otherwise its Taylor series is zero near the origin,
and restriction to each complex line makes it zero everywhere.
Choose such a multi-index $\alpha$.  Its multi-factorial,
$P_\alpha$, and $c_0$ are nonzero.  Comparing (QD.2) for two
scalars, cancel these factors and use injectivity of
$\operatorname{opp}_p$ to get equality.

Take $\kappa=\kappa_{\varnothing}$ from QD15.  For each rational
prime $q$ where both original vertex lattices are standard,
apply QD15 with $Q=\{q\}$.  Its scalar satisfies the same
all-jet identity for the same $G_{x,y}$, so it equals $\kappa$.
Consequently one and the same scalar has
\[
 \mathrm{Mod}_1(\kappa),\qquad
 \mathrm{Unit}_q(\kappa)\quad
  \hbox{at every such }q,\qquad \texttt{JetKap}(x,y,\kappa).
                                                               \tag{QD.21}
\]
This proves the actual conclusion of \texttt{exists\_kappa}.
There is no compactness argument or infinite simultaneous choice
of unrelated scalars.  In particular the selected scalar is
independent of the finite decomposition used to obtain it.

\paragraph{QD17: the literal moved-vacuum dictionary and its eight clauses.}
Choose $L_s=L_s(x,y)$ and $L_t=L_t(x,y)$, and take $\kappa$
from QD16.  These choices prove \texttt{movedVacuum\_dict},
including its geometric and analytic clauses.

The actual moved rational lattices are recovered by intersecting
the finite adelic lattice with the rational vector space in the
unchanged frame.  Duality commutes with this operation: test on
finite generators locally, use the Hermitian inverse of the
unitary element, and intersect the integral local conditions.
They are therefore full self-dual lattices.  Acting on rank-one
Hom generators sends the old Hom lattice to the Hom lattice of
these two moved lattices; linearity and closure give the exact
orbit-lattice inclusion required in the source.  The actual Weil
covariance then gives the moved finite lattice invariance.  Its
finite components commute with real Fock operators and preserve
the fixed real Cartans.  The exact Haar norm is preserved.  Thus
the actual moved vector satisfies the source's full
\texttt{IsVacO}, including its norm condition, on these lattices.

VI2 gives their period-lattice and semicharacter identifications;
VI5 gives the nonzero entire function $G_{x,y}$.  QD16 supplies
its all-conjugate scalar, all local-unit clauses and all jets.
For a top edge the two \texttt{EqAt} conclusions are the two
given vertex conclusions; for a $\chi$ edge they are interchanged,
not replaced by a new frame.  These verify every conclusion of
the literal record.  The source's \texttt{Big} wrapper may take
both extra size thresholds to be zero, because this construction
uses only the supplied valid choice and has no further large-size
condition.  This does not assert that such a choice exists for
arbitrary raw arithmetic data.

\paragraph{QD18: exact size of the split-half quotient.}
Let $\Lambda$ be the actual full self-dual finite Hom lattice
and let $e$ be the orthogonal sum of the split idempotents in
the finite set $P=\operatorname{locPl}(\mathfrak c.G,t)$.
For $k>0$ put $M=\Lambda\cap\bar e\mathcal W_f$ and
$Y=k^{-1}\Lambda\cap\bar e\mathcal W_f$.  Then $Y/M$ is
finite, of order $n>0$, and
\[
 n=\left(\prod_{v\in P_F}|\mathcal O_{F_v}/k\mathcal O_{F_v}|
                                                \right)^{ab},
                           \qquad n\mid k^{f a b}.        \tag{QD.22}
\]
Here $P_F$ contains each supporting base prime once.  The source's
split-place disjointness condition excludes selecting both
conjugate halves, so the passage from $P$ to $P_F$ is injective.

At a split prime $E_v=F_v\times F_v$, projection to either half
identifies the Hom lattice with an integral matrix lattice of
rank $ab$ over $\mathcal O_{F_v}$.  The two original full
lattices are free over that DVR; changing their local bases
does not change the quotient cardinality.  Orthogonal support
idempotents give a direct product over $P_F$.  Multiplication
by $k$ identifies $Y/M$ with $M/kM$, proving the equality.
Finally the local ideal norm product gives
\[
       \prod_{v\text{ finite}}|\mathcal O_{F_v}/k\mathcal O_{F_v}|
            =|\mathcal O_F/k\mathcal O_F|=k^{[F:\mathbb Q]}.
                                                               \tag{QD.23}
\]
Every omitted factor is a positive integer and equals one away
from $k$.  Thus the partial product divides the full product.
This proves the precise exponent $f a b$, including all
ramification multiplicities; replacing it by a guessed residue
depth or by $[E:\mathbb Q]ab$ would lose the source's bound.
For $k=1$, an empty set of places, or $ab=0$, the quotient is
trivial and $n=1$.

\paragraph{QD19: the actual coset indicator has bounded Fourier coefficients.}
Let $d$ be the given \texttt{FinDatum} with
$\texttt{IsCosetDatum}(k,d)$, $k>0$, $d.z=1$ and
$d.P=\operatorname{locPl}(\mathfrak c.G,t)$.
Its vector is a finite Heisenberg combination of the actual
vacuum, with $n$ as in QD18:
\[
       \mathfrak c.S.Z.\operatorname{vec}(d)
                   =\sum_{i=1}^n c_i U(w_i)v_p,
 \quad kw_i\in\Lambda,\quad(w_i)_\infty=0,
 \quad c_i={\zeta_i\over n},\quad\zeta_i^{2k^2}=1.        \tag{QD.24}
\]
Consequently
\[
 k^{f a b}c_i\text{ is an algebraic integer},\qquad
       \sum_i|\rho(c_i)|=1
       \quad\text{for every }\rho:\overline{\mathbb Q}\to\mathbb C.
                                                               \tag{QD.25}
\]

To verify the formula, write the actual coset as $x_0+k\Lambda_e$,
with $kx_0\in\Lambda_e$, and choose representatives $r_i$ of
$k^{-1}\Lambda_{\bar e}/\Lambda_{\bar e}$.  Self-duality makes
the pairing between the two finite half-quotients perfect.
The sum of all its characters is $n$ on the zero coset and
zero elsewhere, by translation by a character taking a nontrivial
value on a nonzero coset.  Thus
\[
 1_{k\Lambda_e}(t)=\frac1n\sum_i
              \psi(\operatorname{sympl}(t,r_i))1_{\Lambda_e}(t).
                                                               \tag{QD.26}
\]
Apply the actual normalized half-Schwartz-model map, which sends
$1_{\Lambda_e}$ to $v_p$, and then translate by $-x_0$ in its
specified convention.  The Heisenberg multiplication law gives
$w_i=-x_0+r_i$ and
$\zeta_i=\psi(\operatorname{symplOH}(-x_0,r_i))$.
Since $2\operatorname{symplOH}=\operatorname{sympl}$ and
$kx_0,kr_i\in\Lambda$, integrality of the full symplectic
pairing gives $\zeta_i^{2k^2}=1$.  This proves (QD.24).
The normalizing map and $d.z=1$ are essential: an additional
Weil translate of the datum is not silently dropped.  QD18
and integrality of roots of unity prove the first part of
(QD.25); every conjugate of $\zeta_i$ has modulus one, so the
second is the sum of $n$ terms $1/n$.

\paragraph{QD20: finite torsion Heisenberg translates have the same period jets.}
For the lattices and scalar of QD17, let $w$ satisfy $kw\in\Lambda$
in the original finite Hom lattice.  There are a root of unity
$\zeta\in\overline{\mathbb Q}$ and a period vector $s$ such that
$ks\in A_{x,y}$ and the actual jets of $U(w)v_p$ at $xu_s,yu_t$
are the jets of
\[
                         \kappa\zeta U(s)G_{x,y},         \tag{QD.27}
\]
with exactly the outside factor $c_0\operatorname{opp}_p(P_\alpha\cdot)$.

First use the actual Weil covariance to move $w$ to
$w'=\operatorname{act}(xu_s,yu_t)w$.  The Hom-lattice inclusion
of QD17 gives $kw'$ in the moved lattice.  Finite rational
approximation supplies a rational Hom vector $M$ with
$u=w'-M_f$ in that lattice.  Then $kM_f=kw'-ku$ is in it too.
The actual lattice-vacuum law gives $U(u)F=\epsilon F$ with
$\epsilon\in\{1,-1\}$.  The two Heisenberg factors therefore
give
$U(w')F=\zeta U(M_f)F$, where
$\zeta=\epsilon\psi(-\operatorname{symplOH}(M_f,u))$.
Since $kM_f,u$ are integral and
$2\operatorname{symplOH}=\operatorname{sympl}$,
$\zeta^{2k}=1$.  This is the finite-rational decomposition in
the source, not a replacement of the adelic vector by an
arbitrary real vector.

Rational quasi-periodicity replaces $M_f$ by its negative real
component.  In the fixed orthonormal coordinates put
\[
 X_\tau=-\,\operatorname{fr}_{T,\tau}^{-1}M_{\infty,\tau}
                         \operatorname{fr}_{S,\tau},
                              \qquad s=\delta X.           \tag{QD.28}
\]
The rational-lattice/period-lattice identification gives
$ks\in A_{x,y}$, with this sign.  Write the real restriction
$F=\mathrm{Gauss}\,\Phi$.  The exact Fock monomial formula and
(QD.2) give $\Phi(Z)=c_0\kappa G_{x,y}(\delta Z)$ after
conjugation in the $\chi$ case.  A real Heisenberg translate
is
\[
 \Phi(Z+X)\exp\{-\pi(2\xi_R\langle Z,X\rangle
                                      +\xi_R\langle X,X\rangle)\}.
                                                               \tag{QD.29}
\]
The source's scaled torus Hermitian form turns (QD.29) exactly
into $U(\delta X)G_{x,y}$; its periods satisfy
$\operatorname{period}(\delta,\tau)
=\delta/\sqrt{2\pi\xi_R}$.  Differentiating gives (QD.27).
For the negative source conjugate the positive calculation and
conjugate back, retaining $\operatorname{opp}_p$ outside the
formula.  All functions differentiated here are entire; the
Heisenberg exponential and translation preserve that property.

\paragraph{QD21: the full quantitative dictionary on the original vector.}
Use QD19's $c_i,w_i$ and apply QD20 to each $w_i$, obtaining
$\zeta_i',s_i$.  Define
\[
             r_i=\operatorname{opp}_{p,\overline{\mathbb Q}}(c_i)
                                         \zeta_i'.         \tag{QD.30}
\]
Then the same $L_s,L_t,\kappa$ of QD17 and these $r_i,s_i$
satisfy every clause of \texttt{DictAt}.  Explicitly,
\[
 \begin{gathered}
  k s_i\in A_{x,y},\qquad
  k^{f a b}r_i\text{ is integral over }\mathbb Z,\qquad
  \sum_i|\rho(r_i)|\le1\quad(\rho:\overline{\mathbb Q}\to\mathbb C),\\
 \operatorname{jetTab}(\operatorname{vec}(d),xu_s,yu_t;\mathrm{mon})
  =c_0\operatorname{opp}_p\left(
        P_\alpha\,\partial^\alpha
          \left[\kappa\sum_i r_iU(s_i)G_{x,y}\right](0)\right),
 \qquad \alpha=\mathrm{mon}.\operatorname{mapDomain}
                                      (\operatorname{fvar}_p).
 \end{gathered}                                            \tag{QD.31}
\]
The all-prime local-unit assertion for $\kappa$ and the two
local lattice assertions are exactly those of QD17.
In fact this uncollected presentation has $r_i=\nu_i/n$ with
$\nu_i^{2k^2}=1$: combine QD19's order $2k^2$ with QD20's
order $2k$.  Thus its coefficients lie in the fixed cyclotomic
field $\mathbb Q(\mu_{2k^2})$, even as the adelic points vary;
this extra observation does not assert such a field for $\kappa$.

Indeed multiplication by a root of unity and the algebraic
involution preserve the integrality in (QD.25).  For a complex
embedding $\rho$, the norm of each $\rho(\zeta_i')$ is one,
so the sum of norms in (QD.31) equals that in (QD.25), evaluated
at $\rho\circ\operatorname{opp}_{p,\overline{\mathbb Q}}$.
Linearity of the given Weil action and the Fock differential
operators applies to (QD.24).  Finite Heisenberg translates
with zero infinite component commute with these real differential
operators.  Apply QD20 term by term.  Move the finite entire
sum through $\partial^\alpha$; the two occurrences of
$\operatorname{opp}_p$ on $c_i$ cancel.  This proves the exact
jet identity.  Omitting the conjugation in (QD.30) would give
the wrong $\chi$-edge coefficients.  No integral, infinite sum
or arbitrary field embedding is interchanged with differentiation.

\paragraph{QD22: the actual Big endpoint and the field boundary.}
QD21 proves \texttt{DictionarySizeBound.dictAt\_big} under
its original literal input hypotheses.  More precisely, it is
valid for every supplied $\mathfrak c.\mathrm{Is}$ in that
source context; QD17 and its finite-translation extension use
both additional size thresholds equal to zero.  Thus the
\texttt{Big} quantifier is satisfied by those same witnesses,
with no growth estimate hidden in the word ``big.''  The
quantitative content is instead the exact torsion level $k$,
denominator $k^{f a b}$, all-conjugate coefficient sum at most
one, and the simultaneous local-unit conclusion.  The finite
presentation even has the positive number of terms $n\mid
k^{f a b}$ obtained in QD18, though the literal \texttt{DictAt}
does not request a separate bound on $n$.

For a fixed finite collection of such dictionaries the finitely
many algebraic coefficients and scalars can of course be put
in one finite extension of the already embedded coefficient
field: adjoin them inside the original algebraic closure, then
use the tower law.  This observation does not produce a single
field for infinitely many arbitrary adelic points, nor does it
put the critical toric ratios in the originally prescribed field.
The surviving \texttt{ToricRatioConjugates}/\texttt{NFRatio}
and original-embedding \texttt{Choice.FieldOK} suppliers remain
separate targets.  The universal uniqueness of the auxiliary
quadratic character is also not asserted.  The real-character
argument of QD8 and CW24's prescribed-character construction
allow the present endpoint to avoid that unproved uniqueness
route while retaining the exact supplied character restrictions.
